\subsection{Divisibility in principal ideal domains}

Let A be a principal ideal domain and let K be its field of fractions (I, p.~116); we will see that the ordered group \(P^{*}\) of principal fractional ideals (VI, p.~6) of K is lattice ordered; more precisely:

PROPOSITION 1. --- Let K be the field of fractions of a principal ideal domain A, and let \((x_{\iota})_{\iota \in I}\) be a family of elements of K having a common denominator \(b \in K^{*}\) (in other words \(bx_{\iota} \in A\) for all \(\iota\) ). Then :

\begin{enumerate}
\def\labelenumi{\alph{enumi})}
\item
  The family \((x_{\iota})\) has a gcd in K.
\item
  Every gcd of \((x_{i})\) can be expressed in the form \(d = \sum_{i}a_{i}x_{i}\) , where the
\end{enumerate}

\(a_{i}\) are elements of A, all but finitely many of which are zero.

Indeed the ideal \(\sum_{i} A b x_{i}\) of \(A\) is principal, and so of the form \(A d'\) . Put

\(d' = bd (d \in \mathbf{K})\) . From the relation \(d' = \sum_{i} a_i b x_i\) , we deduce \(d = \sum_{i} a_i x_i\) , where

\(a_{i} \in A\) . Hence every common divisor of the \(x_{i}\) divides d.~On the other hand, since bd is a common divisor of the \(bx_{i}\) by construction, it follows that d is a common divisor of the \(x_{i}\) .

Remark. --- Prop. 1 applies with no restrictions to an arbitrary family \((x_{i})\) of elements of A (take \(b = 1\) ), and also to any finite family \((x_{i})\) of elements of K (if \(x_{i} = c_{i}b_{i}^{-1}\) with \(c_{i} \in \mathbf{A}\) and \(b_{i} \in \mathbf{A}\) , then take \(b\) to be the product of the \(b_{i}\) ).

COROLLARY. --- Let \((x_{t})\) be an arbitrary family of elements of a principal ideal domain A contained as a subring in an integral domain B, and let d be a gcd of the family \((x_{t})\) in A. Then the family \((x_{t})\) has gcd's in B, and d is one of them.

Indeed \(d\) is a common divisor of the \(x_{i}\) in B. On the other hand the relation \(d = \sum_{i} a_{i} x_{i}\) shows that every common divisor of the \(x_{i}\) in B divides \(d\) .

An important application of this corollary is when \(\mathbf{A} = \mathbf{K}[\mathbf{X}]\) and \(\mathbf{B} = \mathbf{E}[\mathbf{X}]\) , where \(\mathbf{K}\) is a field and \(\mathbf{E}\) is an extension of \(\mathbf{K}\) (IV, p.~12, Cor. 1).

The first assertion of Prop. 1 shows that the ordered group \(P^{*}\) is lattice ordered (VI, p.~10). In particular every finite family of elements of K admits an lcm. We can thus apply the results denoted (DIV) in VI, pp.~10 to 17 to principal ideal domains.

The following result is a consequence of the second assertion of Prop. 1 :

THEOREM 1 (« Bezout's identity »). --- Elements \(x_{\iota}\) ( \(\iota \in I\) ) of a principal ideal domain A are setwise coprime if and only if there exist elements \(a_{\iota}\) ( \(\iota \in I\) ) of A, all but finitely many of them zero, such that \(\sum a_{\iota}x_{\iota} = 1\) .

The necessity of the condition is just Prop. 1. Conversely, if \(\sum_{i} a_i x_i = 1\) then every common divisor of the \(x_i\) divides 1, and so 1 is a gcd of the \(x_i\) .

PROPOSITION 2. --- Let \(a, b, d, m\) and \(p\) be elements of the field of fractions \(K\) of a principal ideal domain \(A\) .

\begin{enumerate}
\def\labelenumi{\alph{enumi})}
\item
  « d is a gcd of a and b » is equivalent to « (d) = (a) + (b) ».
\item
  « m is an lcm of a and b » is equivalent to « m = (a) ∩ (b) ».
\item
  « p is an irreducible element of A » is equivalent to « (p) is a nonzero maximal ideal of A » and to « (p) is a nonzero prime ideal of A ».
\end{enumerate}

We have already proved \(a\) (Prop. 1). Since the common multiples of \(a\) and \(b\) are the elements of \((a) \cap (b)\) , and since, by hypothesis, the ideal \((a) \cap (b)\) is principal, say \((a) \cap (b) = (m)\) , it follows that \(m\) is an lcm of \(a\) and \(b\) , which proves \(b\) ). Finally, to say that \(p \neq 0\) is an irreducible element of \(A\) means by definition (VI, p.~17) that \((p)\) is a maximal element of the family of principal ideals \(\neq A\) of \(A\) , ordered by inclusion; since \(A\) has no ideals other than the principal ideals, this means that \((p)\) is a maximal ideal of \(A\) , whence \(c\) ), by the remark in VI, p.~17.

In a principal ideal domain A, the sum (resp. the intersection) of a finite number of ideals is also called the gcd (resp. lcm) of these ideals.

PROPOSITION 3. --- Let \(a, b, c\) be elements of the field of fractions of a principal ideal domain A, and let \(d\) be a gcd of \(a\) and \(c\) ; then the congruence \(ax \equiv b \pmod{c}\) has a solution \(x_0 \in A\) if and only if \(d\) divides \(b\) ; in this case the elements \(x\) of A which satisfy \(ax \equiv b \pmod{c}\) are precisely those which satisfy \(x \equiv x_0 \pmod{cd^{-1}}\) .

If \(ax \equiv b \pmod{c}\) , with \(x \in \mathbf{A}\) , then there exists \(y \in \mathbf{A}\) such that \(b = ax + cy\) , so \(d\) divides \(b\) . Conversely, if \(d\) divides \(b\) then \(b = ax_0 + cy_0\) with \(x_0, y_0 \in \mathbf{A}\) (Prop. 1), so \(ax_0 \equiv b \pmod{c}\) ; moreover, the relation \(ax \equiv b \pmod{c}\) is then equivalent to \(a(x - x_0) \equiv 0 \pmod{c}\) ; putting \(a = da'\) and \(c = dc'\) , the latter becomes \(a'(x - x_0) \equiv 0 \pmod{c'}\) . But this is equivalent (for \(x \in \mathbf{A}\) ) to \(x - x_0 \equiv 0 \pmod{c'}\) , since \(a'\) and \(c'\) are coprime (VI, p.~14, Prop. 10 (DIV) and VI, p.~15, Cor. 2 to Prop. 11 (DIV)).

PROPOSITION 4. --- Let \((a_i)_{1 \leq i \leq n}\) be a finite family of pairwise coprime elements of the principal ideal domain A. Then the canonical homomorphism (I, p.~110) from \(\mathbf{A}/\left(\prod_{i=1}^{n} a_i\right)\) into the product \(\prod_{i=1}^{n} \mathbf{A}/(a_i)\) is an isomorphism of A-algebras.

This follows from I, p.~110, Prop. 9 and from Prop. 2 a).

The conclusion of Prop. 4 is not valid if it is not assumed that the \(a_{i}\) are pairwise coprime (cf.~VII, p.~24, Prop. 9).

